\documentclass[aspectratio=169]{beamer}
\usepackage[spanish]{babel}
\usepackage[utf8]{inputenc}
\usepackage{amsmath, amssymb, amsthm}
\usepackage{mathtools}
\usepackage{mathrsfs}
\usepackage{tikz}
\usepackage{pgfplots}
\usetikzlibrary{arrows, positioning, automata, calc, patterns, decorations.pathmorphing, matrix, shapes.geometric}
\pgfplotsset{compat=1.18}

% Temática estilo IHES (minimalista, elegante)
\usetheme{default}
\usecolortheme{beaver}
\usefonttheme{serif}

% Configuración de colores
\setbeamercolor{title}{fg=black!90, bg=gray!10}
\setbeamercolor{frametitle}{fg=black!90, bg=gray!10}
\setbeamercolor{block title}{fg=white, bg=black!80}
\setbeamercolor{block body}{fg=black, bg=gray!10}
\setbeamercolor{example block title}{fg=white, bg=blue!50!black!80}
\setbeamercolor{example block body}{fg=black, bg=blue!5}
\setbeamercolor{definition block title}{fg=white, bg=green!50!black!80}
\setbeamercolor{definition block body}{fg=black, bg=green!5}

% Quitar barras de navegación
\setbeamertemplate{navigation symbols}{}

% Estilo de pies de página
\setbeamertemplate{footline}{
  \hfill
  \insertframenumber/\inserttotalframenumber
  \hspace{1em}
  \vspace{0.2em}
}

% Más espacio para imágenes
\setbeamersize{text margin left=1cm, text margin right=1cm}

% Teoremas, definiciones, etc.
\theoremstyle{plain}
\newtheorem{teorema}{Teorema}
\newtheorem{proposicion}{Proposición}
\newtheorem{lema}{Lema}
\newtheorem{corolario}{Corolario}

\theoremstyle{definition}
\newtheorem{definicion}{Definición}
\newtheorem{ejemplo}{Ejemplo}
\newtheorem{observacion}{Observación}

% Atajos
\newcommand{\A}{\mathbf{A}}
\newcommand{\I}{\mathbf{I}}
\newcommand{\PP}{\mathbf{P}}
\newcommand{\p}{\mathbf{p}}
\newcommand{\E}{\mathbb{E}}
\newcommand{\Prob}{\mathbb{P}}
\newcommand{\1}{\mathbf{1}}

\title{Cadenas de Markov en Tiempo Continuo}
\author{César Galindo}
\date{}

\begin{document}

\begin{frame}
  \titlepage
\end{frame}

\begin{frame}{Contenido}
  \tableofcontents
\end{frame}

\section{Propiedades de la Exponencial}

\begin{frame}{Relojes exponenciales independientes}
  \begin{definicion}
    Sean $T_1,\ldots,T_n$ variables aleatorias independientes con distribución exponencial de parámetros $b_1, \ldots, b_n$ respectivamente. Definimos
    \[ T = \min\{T_1, \ldots, T_n\}. \]
    Interpretamos $T_i$ como el momento en que suena el reloj $i$.
  \end{definicion}
  
  \vspace{0.5cm}
  
  \begin{center}
    \begin{tikzpicture}[scale=1.0]
      % Dibujar relojes
      \foreach \i in {1,...,4} {
        \draw (2.2*\i,0) circle (0.5);
        \draw (2.2*\i,0) -- ++(0,0.4);
        \draw (2.2*\i,0) -- ++(0.25,-0.25);
        \node at (2.2*\i,-1) {$T_\i$};
      }
      % Flecha de tiempo
      \draw[->, thick] (0,-1.8) -- (10,-1.8) node[right] {$t$};
      \foreach \x in {0,2.2,4.4,6.6,8.8} {
        \draw (\x,-1.8) -- (\x,-1.7);
      }
      % Marcar tiempos
      \draw[red, thick] (2.8,-1.8) -- (2.8,-1.4) node[above] {$T$};
      \node at (2.8,-2.3) {mínimo};
    \end{tikzpicture}
  \end{center}
\end{frame}

\begin{frame}{El mínimo de exponenciales es exponencial}
  \begin{teorema}
    La variable aleatoria $T = \min\{T_1, \ldots, T_n\}$ tiene distribución exponencial con parámetro $b_1 + \cdots + b_n$.
  \end{teorema}
  
  \begin{proof}
    Calculamos la función de supervivencia:
    \[
    \Prob\{T \geq t\} = \Prob\{T_1 \geq t, \ldots, T_n \geq t\}.
    \]
    Por independencia:
    \[
    \Prob\{T \geq t\} = \prod_{i=1}^n \Prob\{T_i \geq t\} = \prod_{i=1}^n e^{-b_i t} = e^{-(b_1 + \cdots + b_n)t}.
    \]
  \end{proof}
\end{frame}

\begin{frame}{Probabilidad de que un reloj sea el primero}
  \begin{teorema}
    La probabilidad de que el reloj $i$ sea el primero en sonar es
    \[ \Prob\{T_i = T\} = \frac{b_i}{b_1 + \cdots + b_n}. \]
  \end{teorema}

  \vspace{0.3cm}
  Idea: se integra sobre el instante en que suena el reloj $i$ y se exige que todos los demás sigan activos hasta ese tiempo.
\end{frame}

\begin{frame}{Probabilidad de que un reloj sea el primero (demostración)}
  \small
  \begin{proof}
    Condicionamos en el valor de $T_i$:
    \begin{align*}
      \Prob\{T_i = T\} &= \int_0^\infty \Prob\{T_j > t\ \forall j\neq i\} \, b_i e^{-b_i t}\,dt \\
      &= \int_0^\infty e^{-(b_1+\cdots+b_{i-1}+b_{i+1}+\cdots+b_n)t} b_i e^{-b_i t}\,dt \\
      &= b_i \int_0^\infty e^{-(b_1+\cdots+b_n)t}\,dt \\
      &= \frac{b_i}{\sum_j b_j}.
    \end{align*}
  \end{proof}
\end{frame}

\section{Cadenas en Tiempo Continuo}

\begin{frame}{Definición de CTMC}
  \begin{definicion}
    Sea $S$ un espacio de estados finito. Un proceso $\{X_t\}_{t\ge0}$ con valores en $S$ es una \textbf{cadena de Markov en tiempo continuo} si satisface:
    \begin{itemize}
      \item \textbf{Propiedad de Markov:} $\Prob\{X_t = y \mid X_r, 0 \leq r \leq s\} = \Prob\{X_t = y \mid X_s\}$.
      \item \textbf{Homogeneidad temporal:} $\Prob\{X_t = y \mid X_s = x\} = \Prob\{X_{t-s} = y \mid X_0 = x\}$.
    \end{itemize}
  \end{definicion}
\end{frame}

\begin{frame}{Tasas de transición}
  \begin{definicion}
    Para cada $x \neq y$, asignamos una \textbf{tasa de transición} $\alpha(x,y) \geq 0$. La \textbf{tasa total de salida} de $x$ es
    \[ \alpha(x) = \sum_{y \neq x} \alpha(x,y). \]
  \end{definicion}

  \vspace{0.4cm}
  \begin{block}{Interpretación}
    En un intervalo muy pequeño, $\alpha(x,y)$ mide la rapidez con la que el proceso abandona $x$ y salta hacia $y$.
  \end{block}
\end{frame}

\begin{frame}{Probabilidades infinitesimales}
  \begin{block}{Comportamiento en intervalos pequeños}
    \begin{align*}
      \Prob\{X_{t+\Delta t} = x \mid X_t = x\} &= 1 - \alpha(x)\Delta t + o(\Delta t),\\
      \Prob\{X_{t+\Delta t} = y \mid X_t = x\} &= \alpha(x,y)\Delta t + o(\Delta t), \quad y \neq x.
    \end{align*}
  \end{block}
\end{frame}

\begin{frame}{Probabilidades infinitesimales: esquema}
  \begin{center}
    \begin{tikzpicture}[
      state/.style={circle, draw, minimum size=1.0cm, thick},
      ->, >=stealth, thick,
      scale=0.95
    ]
      \node[state] (x) at (0,0) {$x$};
      \node[state] (y1) at (4,1.5) {$y_1$};
      \node[state] (y2) at (4,0) {$y_2$};
      \node[state] (y3) at (4,-1.5) {$y_3$};
      
      \path (x) edge node[above, pos=0.58] {$\alpha(x,y_1)\Delta t$} (y1);
      \path (x) edge node[above, pos=0.58] {$\alpha(x,y_2)\Delta t$} (y2);
      \path (x) edge node[below, pos=0.58] {$\alpha(x,y_3)\Delta t$} (y3);
      \path (x) edge[loop above, min distance=1.5cm] node {$1-\alpha(x)\Delta t$} (x);
    \end{tikzpicture}
  \end{center}
\end{frame}

\begin{frame}{Generador infinitesimal}
  \begin{definicion}
    La \textbf{matriz generadora infinitesimal} $\A$ tiene entradas:
    \[ \A(x,y) = \begin{cases}
      \alpha(x,y) & \text{si } x \neq y,\\
      -\alpha(x) & \text{si } x = y.
    \end{cases} \]
  \end{definicion}

  \begin{observacion}
    Las filas de $\A$ suman cero y las entradas fuera de la diagonal son no negativas.
  \end{observacion}
\end{frame}

\begin{frame}{Forma de la matriz generadora}
  \begin{center}
    \begin{tikzpicture}[scale=0.95]
      \matrix[matrix of math nodes, left delimiter={[}, right delimiter={]}, ampersand replacement=\&] {
        -\alpha(1) \& \alpha(1,2) \& \cdots \& \alpha(1,n) \\
        \alpha(2,1) \& -\alpha(2) \& \cdots \& \alpha(2,n) \\
        \vdots \& \vdots \& \ddots \& \vdots \\
        \alpha(n,1) \& \alpha(n,2) \& \cdots \& -\alpha(n) \\
      };
    \end{tikzpicture}
  \end{center}

  \vspace{0.3cm}
  \[
    \sum_y \A(x,y) = -\alpha(x) + \sum_{y\neq x} \alpha(x,y) = 0.
  \]
\end{frame}

\begin{frame}{Ecuaciones diferenciales}
  \begin{proposicion}
    Las probabilidades $p_x(t) = \Prob\{X_t = x\}$ satisfacen
    \[ p_x'(t) = -\alpha(x)p_x(t) + \sum_{y \neq x} \alpha(y,x)p_y(t). \]
  \end{proposicion}

  \vspace{0.3cm}
  \begin{center}
    \Large
    $\text{flujo que entra} - \text{flujo que sale}$
  \end{center}
\end{frame}

\begin{frame}{Ecuaciones diferenciales: demostración I}
  \begin{proof}
    Condicionando en el estado en el instante $t$:
    \begin{align*}
      p_x(t+\Delta t)
      &= \sum_y \Prob\{X_t=y\}\Prob\{X_{t+\Delta t}=x\mid X_t=y\} \\
      &= p_x(t)\Prob\{X_{t+\Delta t}=x\mid X_t=x\}
         + \sum_{y\neq x} p_y(t)\Prob\{X_{t+\Delta t}=x\mid X_t=y\}.
    \end{align*}
    Usando las probabilidades infinitesimales:
    \[
      p_x(t+\Delta t)
      = p_x(t)(1-\alpha(x)\Delta t)
      + \sum_{y\neq x} p_y(t)\alpha(y,x)\Delta t
      + o(\Delta t).
    \]
  \end{proof}
\end{frame}

\begin{frame}{Ecuaciones diferenciales: demostración II}
  \begin{proof}[Continuación]
    Restamos $p_x(t)$ a ambos lados y dividimos entre $\Delta t$:
    \[
      \frac{p_x(t+\Delta t)-p_x(t)}{\Delta t}
      = -\alpha(x)p_x(t) + \sum_{y\neq x}\alpha(y,x)p_y(t) + o(1).
    \]
    Al tomar el límite cuando $\Delta t\to 0$ obtenemos
    \[
      p_x'(t) = -\alpha(x)p_x(t) + \sum_{y\neq x} \alpha(y,x)p_y(t).
    \]
  \end{proof}
\end{frame}

\begin{frame}{Forma matricial}
  \begin{proposicion}
    Sea $\p(t)$ el vector fila de probabilidades. Entonces
    \[
      \p'(t) = \p(t)\A.
    \]
  \end{proposicion}

  \begin{definicion}
    La matriz de transición $\PP_t$ tiene entradas
    \[
      p_t(x,y) = \Prob\{X_t = y \mid X_0 = x\}.
    \]
    Por tanto,
    \[
      \frac{d}{dt}\PP_t = \PP_t\A,
      \qquad
      \PP_0 = \I.
    \]
  \end{definicion}
\end{frame}

\begin{frame}{Solución matricial}
  \begin{proposicion}
    La solución del sistema es
    \[
      \p(t) = \p(0)e^{t\A}.
    \]
    En particular,
    \[
      \PP_t = e^{t\A}
      = \sum_{n=0}^{\infty} \frac{(t\A)^n}{n!}.
    \]
  \end{proposicion}

  \vspace{0.25cm}
  \begin{center}
    \begin{tikzpicture}[scale=1.05, transform shape]
      \node[draw, rounded corners, thick, fill=gray!5, minimum width=3.8cm, minimum height=1cm] at (0,0) {$\PP_t = e^{t\A}$};
      \draw[->, thick] (-2.3,-0.8) -- (2.3,-0.8) node[midway, below] {evolución temporal};
    \end{tikzpicture}
  \end{center}
\end{frame}

\begin{frame}{Ejemplo: cadena de dos estados}
  \begin{ejemplo}
    Estados $\{0,1\}$, tasas $\alpha(0,1)=1$, $\alpha(1,0)=2$.
    \[ \A = \begin{pmatrix}
      -1 & 1 \\ 2 & -2
    \end{pmatrix}. \]
  \end{ejemplo}

  \begin{center}
    \begin{tikzpicture}[
      state/.style={circle, draw, minimum size=1.2cm},
      ->, >=stealth, thick
    ]
      \node[state] (0) at (0,0) {$0$};
      \node[state] (1) at (4,0) {$1$};
      \path (0) edge[bend left] node[above] {$1$} (1);
      \path (1) edge[bend left] node[below] {$2$} (0);
    \end{tikzpicture}
  \end{center}
\end{frame}

\begin{frame}{Ejemplo: cadena de dos estados (solución)}
  Eigenvalores: $\lambda_1=0$, $\lambda_2=-3$.

  \[
  \PP_t = \begin{pmatrix}
    \frac23+\frac13 e^{-3t} & \frac13-\frac13 e^{-3t} \\[4pt]
    \frac23-\frac23 e^{-3t} & \frac13+\frac23 e^{-3t}
  \end{pmatrix}.
  \]
\end{frame}

\begin{frame}{Ejemplo: cadena de cuatro estados}
  \begin{ejemplo}
    Estados $\{0,1,2,3\}$, generador:
    \[ \A = \begin{pmatrix}
      -3 & 1 & 1 & 1 \\
      2 & -4 & 1 & 1 \\
      2 & 1 & -4 & 1 \\
      1 & 1 & 1 & -3
    \end{pmatrix}. \]
  \end{ejemplo}
\end{frame}

\begin{frame}{Ejemplo: cadena de cuatro estados (grafo)}
  \begin{center}
    \begin{tikzpicture}[
      state/.style={circle, draw, minimum size=11mm, line width=0.9pt, fill=gray!8},
      rate/.style={font=\footnotesize, fill=white, inner sep=1.2pt},
      >={stealth},
      line width=0.95pt,
      scale=0.98,
      transform shape
    ]
      \node[state] (0) at (135:2.25) {$0$};
      \node[state] (1) at (45:2.25) {$1$};
      \node[state] (2) at (-45:2.25) {$2$};
      \node[state] (3) at (-135:2.25) {$3$};

      \draw[->] (0) to[bend left=14] node[rate, above] {$1$} (1);
      \draw[->] (1) to[bend left=14] node[rate, below] {$2$} (0);

      \draw[->] (1) to[bend left=14] node[rate, right] {$1$} (2);
      \draw[->] (2) to[bend left=14] node[rate, left] {$1$} (1);

      \draw[->] (2) to[bend left=14] node[rate, below] {$1$} (3);
      \draw[->] (3) to[bend left=14] node[rate, above] {$1$} (2);

      \draw[->] (3) to[bend left=14] node[rate, left] {$1$} (0);
      \draw[->] (0) to[bend left=14] node[rate, right] {$1$} (3);

      \draw[->] (0) to[out=-18, in=108] node[rate, pos=0.48, above right] {$1$} (2);
      \draw[->] (2) to[out=162, in=-72] node[rate, pos=0.48, below left] {$2$} (0);

      \draw[->] (1) to[out=-162, in=72] node[rate, pos=0.48, above left] {$1$} (3);
      \draw[->] (3) to[out=18, in=-108] node[rate, pos=0.48, below right] {$1$} (1);
    \end{tikzpicture}
  \end{center}

  \vspace{0.15cm}
  \begin{center}
    Grafo simétrico con etiquetas separadas para evitar cruces y encimamientos.
  \end{center}
\end{frame}

\begin{frame}{Distribución estacionaria}
  \begin{teorema}
    Para una cadena irreducible en espacio finito, existe una única distribución $\pi$ tal que $\pi\A = 0$ y
    \[ \lim_{t\to\infty} \PP_t = \mathbf{1}\pi. \]
  \end{teorema}
  
  \vspace{0.3cm}
  
  Para el ejemplo de 4 estados: $\pi = (7/20, 1/5, 1/5, 1/4)$.
\end{frame}

\subsection{Interpretación con tiempos de espera}

\begin{frame}{Tiempo hasta el primer salto}
  \begin{teorema}
    Sea $X_0 = x$ y $T = \inf\{t: X_t \neq x\}$. Entonces $T$ tiene distribución exponencial con parámetro $\alpha(x)$.
  \end{teorema}
  
  \begin{proof}
    Propiedad de pérdida de memoria y $\Prob\{T \leq \Delta t\} = \alpha(x)\Delta t + o(\Delta t)$.
  \end{proof}
  
  \vspace{0.5cm}
  
  \begin{center}
    \begin{tikzpicture}[scale=1.2]
      \draw[->, thick] (0,0) -- (7,0) node[right] {$t$};
      \draw[blue, thick] (0,0.3) -- (2.5,0.3) node[midway, above] {en $x$};
      \fill[red] (2.5,0.3) circle (0.12) node[above] {$T$};
      \draw[blue, thick] (2.5,0.3) -- (5,0.8) node[midway, above] {salta};
      \node at (3.5,-0.8) {$T\sim\text{Exp}(\alpha(x))$};
    \end{tikzpicture}
  \end{center}
\end{frame}

\begin{frame}{Probabilidad del destino del salto}
  \begin{teorema}
    Condicionado a $T$, la probabilidad de saltar a $y \neq x$ es
    \[ \Prob\{X_T = y\} = \frac{\alpha(x,y)}{\alpha(x)}. \]
  \end{teorema}
\end{frame}

\begin{frame}{Probabilidad del destino del salto: demostración}
  \begin{proof}
    Para $y\neq x$, comparamos la probabilidad de que el primer salto ocurra en
    $[t,t+\Delta t]$ y termine en $y$ con la probabilidad total de saltar en ese intervalo:
    \[
      \Prob\{X_T=y\mid T\in[t,t+\Delta t]\}
      = \frac{e^{-\alpha(x)t}\alpha(x,y)\Delta t}{e^{-\alpha(x)t}\alpha(x)\Delta t} + o(1).
    \]
    Cancelando el factor común $e^{-\alpha(x)t}\Delta t$ y dejando $\Delta t\to 0$,
    obtenemos
    \[
      \Prob\{X_T=y\} = \frac{\alpha(x,y)}{\alpha(x)}.
    \]
  \end{proof}
\end{frame}

\begin{frame}{Probabilidad del destino del salto: esquema}
  \begin{center}
    \begin{tikzpicture}[
      state/.style={circle, draw, minimum size=1.25cm, thick, fill=gray!5},
      rate/.style={font=\small, fill=white, inner sep=1.5pt},
      >={stealth},
      line width=1pt,
      scale=1.08,
      transform shape
    ]
      \node[state] (x) at (0,0) {$x$};
      \node[state] (y1) at (4.8,2.2) {$y_1$};
      \node[state] (y2) at (5.2,0) {$y_2$};
      \node[state] (y3) at (4.8,-2.2) {$y_3$};

      \draw[->] (x) -- node[rate, pos=0.56, above, sloped] {$\dfrac{\alpha(x,y_1)}{\alpha(x)}$} (y1);
      \draw[->] (x) -- node[rate, pos=0.56, above] {$\dfrac{\alpha(x,y_2)}{\alpha(x)}$} (y2);
      \draw[->] (x) -- node[rate, pos=0.56, below, sloped] {$\dfrac{\alpha(x,y_3)}{\alpha(x)}$} (y3);
    \end{tikzpicture}
  \end{center}
\end{frame}

\begin{frame}{Interpretación de los relojes}
  \begin{observacion}[Relojes independientes]
    En cada estado $y \neq x$ colocamos un reloj independiente que suena con tiempo exponencial de tasa $\alpha(x,y)$. La cadena permanece en $x$ hasta que suena el primer reloj y entonces salta al estado correspondiente.
  \end{observacion}

  \vspace{0.35cm}
  \begin{block}{Idea central}
    El \emph{mínimo} de esos relojes determina el tiempo de espera en $x$, y el reloj ganador determina el siguiente estado.
  \end{block}
\end{frame}

\begin{frame}{Relojes independientes: representación gráfica}
  \begin{center}
    \begin{tikzpicture}[
      state/.style={circle, draw, minimum size=1.15cm, thick, fill=gray!5},
      every node/.style={font=\small},
      >={stealth},
      scale=1.0,
      transform shape
    ]
      \node[state] (x) at (0,0) {$x$};
      \node[state] (y1) at (4.7,2.2) {$y_1$};
      \node[state] (y2) at (4.7,0) {$y_2$};
      \node[state] (y3) at (4.7,-2.2) {$y_3$};

      \draw[dashed, ->, thick] (x) -- (y1) node[midway, above] {$\alpha(x,y_1)$};
      \draw[dashed, ->, thick] (x) -- (y2) node[midway, above] {$\alpha(x,y_2)$};
      \draw[dashed, ->, thick] (x) -- (y3) node[midway, below] {$\alpha(x,y_3)$};

      \node[align=center] at (2.3,3.05) {Relojes exponenciales\\independientes};
    \end{tikzpicture}
  \end{center}
\end{frame}

\section{Galería de Imágenes}

\begin{frame}{Visualización: convergencia estacionaria}
  \begin{center}
    \begin{tikzpicture}
      \begin{axis}[
        width=0.9\textwidth,
        height=0.58\textheight,
        xlabel={$t$},
        ylabel={$p_t(0,0)$},
        domain=0:5,
        samples=100,
        grid=major,
        title={Convergencia a distribución estacionaria}
      ]
      \addplot[blue, thick] {2/3 + (1/3)*exp(-3*x)};
      \addplot[red, dashed, thick] coordinates {(0,2/3) (5,2/3)};
      \legend{{$p_t(0,0)$}, {$\pi(0)=2/3$}}
      \end{axis}
    \end{tikzpicture}
  \end{center}
\end{frame}

\begin{frame}{Visualización: crecimiento del proceso de Yule}
  \begin{center}
    \begin{tikzpicture}
      \begin{axis}[
        width=0.9\textwidth,
        height=0.58\textheight,
        xlabel={$t$},
        ylabel={$\E[X_t]$},
        domain=0:3,
        samples=100,
        grid=major,
        title={Crecimiento exponencial}
      ]
      \addplot[blue, thick] {exp(x)};
      \end{axis}
    \end{tikzpicture}
  \end{center}
\end{frame}

\begin{frame}{Visualización: distribución estacionaria M/M/$\infty$}
  \begin{center}
    \begin{tikzpicture}
      \begin{axis}[
        width=0.9\textwidth,
        height=0.58\textheight,
        xlabel={$n$},
        ylabel={$\pi(n)$},
        domain=0:10,
        samples=50,
        grid=major,
        title={Distribución estacionaria M/M/$\infty$}
      ]
      \addplot[blue, thick] {exp(-2)*2^x/x!};
      \end{axis}
    \end{tikzpicture}
  \end{center}
\end{frame}

\begin{frame}{Visualización: probabilidades del proceso de Yule}
  \begin{center}
    \begin{tikzpicture}
      \begin{axis}[
        width=0.9\textwidth,
        height=0.58\textheight,
        xlabel={$t$},
        ylabel={$P_n(t)$},
        domain=0:3,
        samples=100,
        grid=major,
        title={Proceso de Yule $P_n(t)$}
      ]
      \addplot[blue, thick] {exp(-x)};
      \addplot[red, thick] {exp(-x)*(1-exp(-x))};
      \addplot[green, thick] {exp(-x)*(1-exp(-x))^2};
      \legend{{$n=1$}, {$n=2$}, {$n=3$}}
      \end{axis}
    \end{tikzpicture}
  \end{center}
\end{frame}

\begin{frame}{Estructuras de cadenas de Markov I}
  \begin{center}
    \begin{tikzpicture}[
      state/.style={circle, draw, minimum size=1cm, thick, fill=gray!5},
      >={stealth},
      ->, thick,
      scale=1.15,
      transform shape
    ]
      \node[state] (0) at (0,0) {$0$};
      \node[state] (1) at (2.4,0) {$1$};
      \path (0) edge[bend left=18] node[above] {$\lambda$} (1);
      \path (1) edge[bend left=18] node[below] {$\mu$} (0);
    \end{tikzpicture}

    \vspace{0.25cm}
    Cadena de 2 estados
  \end{center}
\end{frame}

\begin{frame}{Estructuras de cadenas de Markov I: cadena de 4 estados}
  \begin{center}
    \begin{tikzpicture}[
      state/.style={circle, draw, minimum size=1.0cm, thick, fill=gray!5},
      every node/.style={font=\small},
      >={stealth},
      ->, thick,
      scale=1.02,
      transform shape
    ]
      \node[state] (0) at (0,3.0) {$0$};
      \node[state] (1) at (4.0,3.0) {$1$};
      \node[state] (2) at (4.0,-0.2) {$2$};
      \node[state] (3) at (0,-0.2) {$3$};

      \path (0) edge[bend left=14] node[above] {$\alpha$} (1);
      \path (1) edge[bend left=14] node[below] {$\beta$} (0);
      \path (1) edge[bend left=14] node[right] {$\alpha$} (2);
      \path (2) edge[bend left=14] node[left] {$\alpha$} (1);
      \path (2) edge[bend left=14] node[below] {$\alpha$} (3);
      \path (3) edge[bend left=14] node[above] {$\alpha$} (2);
      \path (3) edge[bend left=14] node[left] {$\alpha$} (0);
      \path (0) edge[bend left=14] node[right] {$\alpha$} (3);
      \path (0) edge[bend left=26] node[sloped, above] {$\alpha$} (2);
      \path (2) edge[bend left=26] node[sloped, below] {$\beta$} (0);
      \path (1) edge[bend left=26] node[sloped, above] {$\alpha$} (3);
      \path (3) edge[bend left=26] node[sloped, below] {$\alpha$} (1);
    \end{tikzpicture}

    \vspace{0.25cm}
    Cadena de 4 estados con mayor separación entre nodos y etiquetas.
  \end{center}
\end{frame}

\begin{frame}{Estructuras de cadenas de Markov II}
  \begin{center}
    \begin{tikzpicture}[
      state/.style={circle, draw, minimum size=0.8cm},
      ->, >=stealth, thick, scale=1.05
    ]
      \node[state] (0) at (0,0) {$0$};
      \node[state] (1) at (2.3,0) {$1$};
      \node[state] (2) at (4.6,0) {$2$};
      \node[state] (3) at (6.9,0) {$3$};
      \node at (8.3,0) {$\cdots$};
      \path (0) edge node[above] {$\lambda$} (1);
      \path (1) edge node[above] {$\lambda$} (2);
      \path (2) edge node[above] {$\lambda$} (3);
      \path (1) edge[bend left] node[below] {$\mu$} (0);
      \path (2) edge[bend left] node[below] {$\mu$} (1);
      \path (3) edge[bend left] node[below] {$\mu$} (2);
    \end{tikzpicture}
  \end{center}

  \begin{center}
    Proceso de nacimiento--muerte
  \end{center}
\end{frame}

\begin{frame}{Explosión en tiempo finito: esquema}
  \begin{center}
    \begin{tikzpicture}[scale=1.35]
      \draw[->, thick] (0,0) -- (4.2,0) node[right] {$t$};
      \draw[->, thick] (0,0) -- (0,3.2) node[above] {$X_t$};
      \draw[blue, very thick] plot[smooth] coordinates {(0,0.08) (0.6,0.14) (1.2,0.38) (1.8,0.95) (2.2,1.55) (2.8,2.15) (3.15,2.75)};
      \draw[red, dashed, very thick] (3.2,0) -- (3.2,2.95);
      \node[below] at (3.2,0) {$Y_\infty$};
    \end{tikzpicture}
  \end{center}

  \begin{center}
    El proceso acumula infinitos saltos antes de un tiempo aleatorio finito.
  \end{center}
\end{frame}

\begin{frame}{Estructuras de cadenas de Markov III}
  \begin{center}
    \begin{tikzpicture}[scale=1.15]
      \draw (0,0) rectangle (3,3);
      \node at (1.5,1.5) {$\A$};
      \draw[red, thick] (1,0) -- (1,3);
      \draw[red, thick] (0,1) -- (3,1);
      \node at (0.55,2.45) {$\tilde{\A}$};
    \end{tikzpicture}
  \end{center}

  \vspace{0.4cm}
  \begin{center}
    Matriz reducida para tiempos medios de paso
  \end{center}
\end{frame}

\section{Tiempos medios de paso}

\begin{frame}{Definición}
  \begin{definicion}
    Sea $z$ un estado fijo. Definimos el \textbf{tiempo medio de paso} a $z$ comenzando en $x$ como
    \[ b(x) = \E(Y \mid X_0 = x), \quad \text{donde } Y = \inf\{t: X_t = z\}. \]
    Claramente $b(z) = 0$.
  \end{definicion}
  
  \vspace{0.5cm}
  
  \begin{center}
    \begin{tikzpicture}[scale=1.1]
      \draw[->, thick] (0,0) -- (8,0) node[right] {$t$};
      \draw[blue, thick] (0,1) -- (2,1) node[midway, above] {$x$};
      \fill[red] (2,1) circle (0.1);
      \draw[blue, thick] (2,1) -- (4,0.5) -- (5.5,0) -- (6.5,0.5);
      \fill[green] (6.5,0.5) circle (0.1) node[above] {$z$};
      \draw[<->, thick] (0,-0.8) -- (6.5,-0.8) node[midway, below] {$Y$};
    \end{tikzpicture}
  \end{center}
\end{frame}

\begin{frame}{Ecuación para los tiempos medios de paso}
  \begin{teorema}
    Para $x \neq z$, la función $b(x)$ satisface
    \[ \alpha(x)b(x) = 1 + \sum_{y \neq x, z} \alpha(x,y) b(y). \]
  \end{teorema}
  
  \begin{proof}
    Condicionando en el primer salto:
    \[ b(x) = \E(T) + \sum_y \Prob\{X_T=y\} b(y) = \frac{1}{\alpha(x)} + \sum_{y\neq x} \frac{\alpha(x,y)}{\alpha(x)} b(y). \]
    Usando $b(z)=0$ y multiplicando por $\alpha(x)$ obtenemos la ecuación.
  \end{proof}
\end{frame}

\begin{frame}{Forma matricial}
  \begin{corolario}
    Sea $\tilde{\A}$ la matriz obtenida de $\A$ eliminando la fila y columna correspondientes al estado $z$.
    Sea $\bar{b}$ el vector de tiempos medios para $x \neq z$. Entonces
    \[
      \tilde{\A}\,\bar{b} = -\bar{1}.
    \]
    En consecuencia,
    \[
      \bar{b} = [-\tilde{\A}]^{-1}\bar{1}.
    \]
  \end{corolario}

  \vspace{0.2cm}
  \begin{center}
    \begin{tikzpicture}[scale=0.58, transform shape]
      \draw[line width=0.9pt] (0,0) rectangle (4,4);
      \node at (2,2) {$\A$};
      \draw[red, line width=1pt] (1.5,0) -- (1.5,4);
      \draw[red, line width=1pt] (0,1.5) -- (4,1.5);
      \node[fill=white, inner sep=1pt] at (0.75,3.2) {$\tilde{\A}$};
    \end{tikzpicture}
  \end{center}
\end{frame}

\begin{frame}[t]{Ejemplo numérico}
  Para la cadena de 4 estados con $z=3$:
  \[ \tilde{\A} = \begin{pmatrix}
    -3 & 1 & 1 \\
    2 & -4 & 1 \\
    2 & 1 & -4
  \end{pmatrix}. \]
  
  Resolvemos $-\tilde{\A}\bar{b} = \bar{1}$:
  \[ \begin{cases}
    3b_0 - b_1 - b_2 = 1 \\
    -2b_0 + 4b_1 - b_2 = 1 \\
    -2b_0 - b_1 + 4b_2 = 1
  \end{cases} \]
  
  Solución: $b_0 = b_1 = b_2 = 1$.
  
  \vspace{0.3cm}
  
  El tiempo esperado del estado 0 al estado 3 es 1.
\end{frame}

\section{Procesos de Nacimiento y Muerte}

\begin{frame}{Definición}
  \begin{definicion}
    Un \textbf{proceso de nacimiento y muerte} es una cadena de Markov en tiempo continuo con espacio de estados $\{0,1,2,\ldots\}$ donde los únicos cambios posibles son:
    \begin{itemize}
      \item $n \to n+1$ (nacimiento) con tasa $\lambda_n \geq 0$
      \item $n \to n-1$ (muerte) con tasa $\mu_n \geq 0$, $\mu_0 = 0$
    \end{itemize}
  \end{definicion}
  
  \vspace{0.5cm}
  
  \begin{center}
    \begin{tikzpicture}[
      state/.style={circle, draw, minimum size=0.9cm},
      ->, >=stealth, thick
    ]
      \node[state] (0) at (0,0) {$0$};
      \node[state] (1) at (2.5,0) {$1$};
      \node[state] (2) at (5,0) {$2$};
      \node[state] (3) at (7.5,0) {$3$};
      \node at (9,0) {$\cdots$};
      
      \path (0) edge node[above] {$\lambda_0$} (1);
      \path (1) edge node[above] {$\lambda_1$} (2);
      \path (2) edge node[above] {$\lambda_2$} (3);
      \path (1) edge[bend left] node[below] {$\mu_1$} (0);
      \path (2) edge[bend left] node[below] {$\mu_2$} (1);
      \path (3) edge[bend left] node[below] {$\mu_3$} (2);
    \end{tikzpicture}
  \end{center}
\end{frame}

\begin{frame}{Probabilidades infinitesimales}
  \begin{align*}
    \Prob\{X_{t+\Delta t} = n \mid X_t = n\} &= 1 - (\lambda_n + \mu_n)\Delta t + o(\Delta t),\\
    \Prob\{X_{t+\Delta t} = n+1 \mid X_t = n\} &= \lambda_n \Delta t + o(\Delta t),\\
    \Prob\{X_{t+\Delta t} = n-1 \mid X_t = n\} &= \mu_n \Delta t + o(\Delta t).
  \end{align*}
\end{frame}

\begin{frame}{Probabilidades infinitesimales: esquema del salto}
  \begin{center}
    \begin{tikzpicture}[
      scale=1.08,
      transform shape,
      state/.style={circle, draw, minimum size=1.15cm, thick, fill=gray!5},
      >={stealth},
      ->, thick,
      every node/.style={font=\small}
    ]
      \node[state] (n) at (0,0) {$n$};
      \node[state] (np1) at (4.2,1.8) {$n+1$};
      \node[state] (nm1) at (4.2,-1.8) {$n-1$};
      
      \path (n) edge node[above] {$\lambda_n\Delta t$} (np1);
      \path (n) edge node[below] {$\mu_n\Delta t$} (nm1);
      \path (n) edge[loop above, min distance=1.8cm] node {$1-(\lambda_n+\mu_n)\Delta t$} (n);
    \end{tikzpicture}
  \end{center}
\end{frame}

\begin{frame}{Ecuaciones diferenciales}
  \begin{proposicion}
    Las probabilidades $P_n(t) = \Prob\{X_t = n\}$ satisfacen:
    \[ P_n'(t) = \mu_{n+1}P_{n+1}(t) + \lambda_{n-1}P_{n-1}(t) - (\mu_n + \lambda_n)P_n(t). \]
    con $\lambda_{-1}=0$, $\mu_0=0$.
  \end{proposicion}
  
  \vspace{0.5cm}
  
  \begin{center}
    \begin{tikzpicture}[scale=1.1]
      \node at (0,1.2) {$\mu_{n+1}P_{n+1}$};
      \node at (0,0) {$\lambda_{n-1}P_{n-1}$};
      \node at (0,-1.2) {$-(\mu_n+\lambda_n)P_n$};
      \draw[->, thick] (1,1.2) -- (3,0.6) node[midway, above] {flujo};
      \draw[->, thick] (1,0) -- (3,-0.6) node[midway, above] {flujo};
      \node[circle, draw, minimum size=1cm, thick] at (4,0) {$P_n'$};
    \end{tikzpicture}
  \end{center}
\end{frame}

\subsection{Ejemplos}

\begin{frame}{Ejemplo 1: Proceso de Poisson}
  \begin{ejemplo}
    $\lambda_n = \lambda$, $\mu_n = 0$ para todo $n$.
  \end{ejemplo}
  
  \vspace{0.3cm}
  
  \begin{center}
    \begin{tikzpicture}[
      state/.style={circle, draw, minimum size=0.9cm},
      ->, >=stealth, thick
    ]
      \node[state] (0) at (0,0) {$0$};
      \node[state] (1) at (2.5,0) {$1$};
      \node[state] (2) at (5,0) {$2$};
      \node[state] (3) at (7.5,0) {$3$};
      \node at (9,0) {$\cdots$};
      
      \path (0) edge node[above] {$\lambda$} (1);
      \path (1) edge node[above] {$\lambda$} (2);
      \path (2) edge node[above] {$\lambda$} (3);
    \end{tikzpicture}
  \end{center}
  
  \vspace{0.3cm}
  
  Ecuaciones: $P_n'(t) = \lambda P_{n-1}(t) - \lambda P_n(t)$.
  
  Solución: $P_n(t) = e^{-\lambda t}\frac{(\lambda t)^n}{n!}$.
\end{frame}

\begin{frame}{Ejemplo 2: Cola M/M/1}
  \begin{ejemplo}
    Un servidor, llegadas Poisson con tasa $\lambda$, servicios exponenciales con tasa $\mu$.
    \[ \lambda_n = \lambda,\quad \mu_n = \mu\ (n\ge 1),\ \mu_0=0. \]
  \end{ejemplo}
  
  \vspace{0.3cm}
  
  \begin{center}
    \begin{tikzpicture}[
      state/.style={circle, draw, minimum size=0.9cm},
      ->, >=stealth, thick
    ]
      \node[state] (0) at (0,0) {$0$};
      \node[state] (1) at (2.5,0) {$1$};
      \node[state] (2) at (5,0) {$2$};
      \node[state] (3) at (7.5,0) {$3$};
      \node at (9,0) {$\cdots$};
      
      \path (0) edge node[above] {$\lambda$} (1);
      \path (1) edge node[above] {$\lambda$} (2);
      \path (2) edge node[above] {$\lambda$} (3);
      \path (1) edge[bend left] node[below] {$\mu$} (0);
      \path (2) edge[bend left] node[below] {$\mu$} (1);
      \path (3) edge[bend left] node[below] {$\mu$} (2);
    \end{tikzpicture}
  \end{center}
\end{frame}

\begin{frame}{Ejemplo 3: Cola M/M/k}
  \begin{ejemplo}
    $k$ servidores, cada uno sirve a tasa $\mu$.
    \[ \lambda_n = \lambda,\quad \mu_n = \min(n,k)\mu. \]
  \end{ejemplo}
  
  \vspace{0.3cm}
  
  \begin{center}
    \begin{tikzpicture}[
      state/.style={circle, draw, minimum size=0.8cm},
      ->, >=stealth, thick, scale=0.9
    ]
      \node[state] (0) at (0,0) {$0$};
      \node[state] (1) at (1.8,0) {$1$};
      \node[state] (2) at (3.6,0) {$2$};
      \node[state] (k1) at (6,0) {$k-1$};
      \node[state] (k) at (7.8,0) {$k$};
      \node[state] (kp1) at (9.6,0) {$k+1$};
      \node at (11,0) {$\cdots$};
      
      \path (0) edge node[above] {$\lambda$} (1);
      \path (1) edge node[above] {$\lambda$} (2);
      \path (2) edge node[above, font=\small] {$\lambda$} (k1);
      \path (k1) edge node[above, font=\small] {$\lambda$} (k);
      \path (k) edge node[above, font=\small] {$\lambda$} (kp1);
      
      \path (1) edge[bend left] node[below, font=\small] {$\mu$} (0);
      \path (2) edge[bend left] node[below, font=\small] {$2\mu$} (1);
      \path (k1) edge[bend left] node[below, font=\small] {$(k-1)\mu$} (2);
      \path (k) edge[bend left] node[below, font=\small] {$k\mu$} (k1);
      \path (kp1) edge[bend left] node[below, font=\small] {$k\mu$} (k);
    \end{tikzpicture}
  \end{center}
\end{frame}

\begin{frame}{Ejemplo 4: Cola M/M/$\infty$}
  \begin{ejemplo}
    Infinitos servidores: cada persona en cola es servida inmediatamente.
    \[ \lambda_n = \lambda,\quad \mu_n = n\mu. \]
  \end{ejemplo}
  
  \vspace{0.3cm}
  
  \begin{center}
    \begin{tikzpicture}[
      state/.style={circle, draw, minimum size=0.9cm},
      ->, >=stealth, thick
    ]
      \node[state] (0) at (0,0) {$0$};
      \node[state] (1) at (2.5,0) {$1$};
      \node[state] (2) at (5,0) {$2$};
      \node[state] (3) at (7.5,0) {$3$};
      \node at (9,0) {$\cdots$};
      
      \path (0) edge node[above] {$\lambda$} (1);
      \path (1) edge node[above] {$\lambda$} (2);
      \path (2) edge node[above] {$\lambda$} (3);
      \path (1) edge[bend left] node[below] {$\mu$} (0);
      \path (2) edge[bend left] node[below] {$2\mu$} (1);
      \path (3) edge[bend left] node[below] {$3\mu$} (2);
    \end{tikzpicture}
  \end{center}
\end{frame}

\begin{frame}{Ejemplo 5: Modelo de población}
  \begin{ejemplo}
    Cada individuo produce descendencia a tasa $\lambda$ y muere a tasa $\mu$, independientemente.
    \[ \lambda_n = n\lambda,\quad \mu_n = n\mu. \]
  \end{ejemplo}
  
  \vspace{0.3cm}
  
  \begin{center}
    \begin{tikzpicture}
      \begin{axis}[
        width=0.8\textwidth,
        height=0.4\textheight,
        xlabel={$t$},
        ylabel={$\E[X_t]$},
        domain=0:3,
        samples=100,
        grid=major,
        title={Crecimiento/Decrecimiento poblacional}
      ]
      \addplot[blue, thick] {exp(x)};
      \addplot[red, thick] {exp(-x)};
      \legend{{$\lambda>\mu$}, {$\lambda<\mu$}}
      \end{axis}
    \end{tikzpicture}
  \end{center}
\end{frame}

\begin{frame}{Ejemplo 6: Crecimiento rápido}
  \begin{ejemplo}
    $\lambda_n = n^2\lambda$, $\mu_n = 0$.
  \end{ejemplo}
  
  \vspace{0.3cm}
  
  \begin{center}
    \begin{tikzpicture}
      \begin{axis}[
        width=0.8\textwidth,
        height=0.4\textheight,
        xlabel={$t$},
        ylabel={$X_t$},
        domain=0:3.5,
        samples=100,
        grid=major,
        title={Explosión en tiempo finito}
      ]
      \addplot[blue, thick] {2/(1+exp(-8*(x-2)))};
      \addplot[red, dashed, thick] coordinates {(3.2,0) (3.2,2)};
      \end{axis}
    \end{tikzpicture}
  \end{center}
  
  \begin{center}
    $\displaystyle Y_\infty = \sum_{i=1}^\infty \frac{1}{i^2\lambda} = \frac{\pi^2}{6\lambda} < \infty$
  \end{center}
\end{frame}

\subsection{Clasificación}

\begin{frame}{Criterio de transiencia}
  \begin{teorema}
    Un proceso de nacimiento y muerte irreducible es \textbf{transitorio} si y solo si
    \[ \sum_{n=1}^{\infty} \frac{\mu_1 \mu_2 \cdots \mu_n}{\lambda_1 \lambda_2 \cdots \lambda_n} < \infty. \]
  \end{teorema}
  
  \begin{center}
    \begin{tikzpicture}
      \begin{axis}[
        width=0.7\textwidth,
        height=0.3\textheight,
        xlabel={$n$},
        ylabel={término $n$-ésimo},
        domain=1:10,
        samples=50,
        grid=major
      ]
      \addplot[blue, thick] {1/x};
      \end{axis}
    \end{tikzpicture}
  \end{center}
  
  \begin{center}
    Convergencia $\Rightarrow$ transiencia
  \end{center}
\end{frame}

\begin{frame}{Criterio de recurrencia positiva}
  \begin{teorema}
    Un proceso de nacimiento y muerte irreducible es \textbf{positivamente recurrente} si y solo si
    \[ \sum_{n=0}^{\infty} \frac{\lambda_0 \lambda_1 \cdots \lambda_{n-1}}{\mu_1 \mu_2 \cdots \mu_n} < \infty. \]
    En tal caso, la distribución estacionaria es
    \[ \pi(n) = \frac{\lambda_0 \lambda_1 \cdots \lambda_{n-1}}{\mu_1 \mu_2 \cdots \mu_n} \cdot \frac{1}{C}, \]
    donde $C$ es la suma de la serie.
  \end{teorema}
  
  \begin{center}
    \begin{tikzpicture}
      \begin{axis}[
        width=0.7\textwidth,
        height=0.3\textheight,
        xlabel={$n$},
        ylabel={$\pi(n)$},
        domain=0:10,
        samples=50,
        grid=major
      ]
      \addplot[blue, thick] {exp(-x)};
      \end{axis}
    \end{tikzpicture}
  \end{center}
\end{frame}

\begin{frame}{Clasificación de colas M/M}
  \begin{center}
    \begin{tabular}{|c|c|c|c|}
      \hline
      Modelo & Transitoria & Positivamente recurrente & Distribución $\pi(n)$ \\
      \hline
      M/M/1 & $\mu < \lambda$ & $\lambda < \mu$ & $(1-\rho)\rho^n$, $\rho=\lambda/\mu$ \\
      \hline
      M/M/k & $k\mu < \lambda$ & $\lambda < k\mu$ & fórmula complicada \\
      \hline
      M/M/$\infty$ & nunca & siempre & $e^{-\lambda/\mu}\frac{(\lambda/\mu)^n}{n!}$ \\
      \hline
    \end{tabular}
  \end{center}
  
  \vspace{0.5cm}
  Para M/M/1: $\displaystyle \E[X] = \frac{\lambda}{\mu-\lambda}$.
  
  Para M/M/$\infty$: $\E[X] = \frac{\lambda}{\mu}$.
\end{frame}

\section{Procesos de Nacimiento Puro}

\begin{frame}{Proceso de Yule}
  \begin{definicion}
    Proceso de Yule: $\lambda_n = n\lambda$, $\mu_n = 0$, $X_0 = 1$.
  \end{definicion}
  
  \begin{teorema}
    Las probabilidades son
    \[ P_n(t) = e^{-\lambda t}(1 - e^{-\lambda t})^{n-1}, \quad n \ge 1. \]
  \end{teorema}
  
  \begin{center}
    \begin{tikzpicture}
      \begin{axis}[
        width=0.8\textwidth,
        height=0.3\textheight,
        xlabel={$t$},
        ylabel={$P_n(t)$},
        domain=0:3,
        samples=100,
        legend pos=north east,
        grid=major
      ]
      \addplot[blue, thick] {exp(-x)};
      \addplot[red, thick] {exp(-x)*(1-exp(-x))};
      \addplot[green, thick] {exp(-x)*(1-exp(-x))^2};
      \legend{{$n=1$}, {$n=2$}, {$n=3$}}
      \end{axis}
    \end{tikzpicture}
  \end{center}
\end{frame}

\begin{frame}{Propiedades del proceso de Yule}
  \begin{proposicion}
    Para el proceso de Yule:
    \begin{itemize}
      \item Esperanza: $\E[X_t] = e^{\lambda t}$
      \item Tiempo hasta $n$ individuos: $\displaystyle\E[Y_n] = \sum_{i=1}^{n-1} \frac{1}{i\lambda} \sim \frac{\ln n}{\lambda}$
    \end{itemize}
  \end{proposicion}
  
  \begin{center}
    \begin{tikzpicture}
      \begin{axis}[
        width=0.8\textwidth,
        height=0.3\textheight,
        xlabel={$t$},
        ylabel={$\E[X_t]$},
        domain=0:3,
        samples=100,
        grid=major
      ]
      \addplot[blue, thick] {exp(x)};
      \end{axis}
    \end{tikzpicture}
  \end{center}
\end{frame}

\begin{frame}{Explosión en tiempo finito}
  \begin{ejemplo}[Crecimiento rápido]
    $\lambda_n = n^2\lambda$, $\mu_n = 0$, $X_0 = 1$.
  \end{ejemplo}
  
  El tiempo para alcanzar $n$ es $Y_n = \sum_{i=1}^{n-1} T_i$ con $T_i \sim \text{Exp}(i^2\lambda)$. Entonces
  \[ \E[Y_\infty] = \sum_{i=1}^\infty \frac{1}{i^2\lambda} = \frac{\pi^2}{6\lambda} < \infty. \]
  
  \begin{center}
    \begin{tikzpicture}
      \begin{axis}[
        width=0.8\textwidth,
        height=0.3\textheight,
        xlabel={$t$},
        ylabel={$X_t$},
        domain=0:3.5,
        samples=100,
        grid=major
      ]
      \addplot[blue, thick] {2/(1+exp(-8*(x-2)))};
      \addplot[red, dashed, thick] coordinates {(3.2,0) (3.2,2)};
      \end{axis}
    \end{tikzpicture}
  \end{center}
  
  \begin{center}
    ¡La población se vuelve infinita en tiempo finito con probabilidad 1!
  \end{center}
\end{frame}

\section{Aplicaciones en Inteligencia Artificial}

\begin{frame}[t]{Modelos de Lenguaje (LLMs): idea básica}
  \begin{block}{Generación de texto}
    Los modelos como GPT pueden verse, de manera simplificada, como sistemas que pasan de un contexto a otro:
    \begin{itemize}
      \item estados = secuencias o contextos de tokens,
      \item transiciones = probabilidad o intensidad del siguiente token,
      \item evolución = actualización sucesiva del contexto.
    \end{itemize}
  \end{block}
\end{frame}

\begin{frame}[t]{Modelos de Lenguaje (LLMs): grafo de transición}
  \begin{center}
    \begin{tikzpicture}[
      state/.style={rectangle, draw, minimum width=1.7cm, minimum height=0.85cm, rounded corners, thick, fill=gray!5},
      >={stealth},
      ->, thick,
      every node/.style={font=\small},
      scale=1.05,
      transform shape
    ]
      \node[state] (s1) at (0,1.6) {``el''};
      \node[state] (s2) at (3.4,2.4) {``gato''};
      \node[state] (s3) at (3.4,1.2) {``perro''};
      \node[state] (s4) at (3.4,0) {``coche''};
      
      \path (s1) edge node[above, sloped] {$\alpha$} (s2);
      \path (s1) edge node[above] {$\beta$} (s3);
      \path (s1) edge node[below, sloped] {$\gamma$} (s4);
    \end{tikzpicture}
  \end{center}
\end{frame}

\begin{frame}[t]{Modelos de Lenguaje (LLMs): atención y tiempos}
  \begin{block}{Interpretación cualitativa}
    \begin{itemize}
      \item Tiempos de atención $\leftrightarrow$ tiempos de salto.
      \item Mecanismo de atención $\leftrightarrow$ pesos o tasas de transición.
      \item Estabilidad de la generación $\leftrightarrow$ comportamiento estacionario.
    \end{itemize}
  \end{block}

  \begin{center}
    \begin{tikzpicture}[scale=1.05, transform shape]
      \draw[->, thick] (0,0) -- (4.4,0) node[right] {tiempo};
      \foreach \x/\t in {0.6/a, 1.6/b, 2.7/c, 3.8/d} {
        \draw[blue, thick] (\x,0.1) -- (\x,0.55);
        \node at (\x,0.8) {\t};
      }
    \end{tikzpicture}
  \end{center}
\end{frame}

\begin{frame}[t]{Procesos de Decisión Markovianos (MDP)}
  \begin{block}{Aprendizaje por refuerzo}
    \begin{itemize}
      \item Q-learning: estima recompensas y decisiones futuras.
      \item Políticas óptimas: seleccionan acciones que modifican la dinámica.
      \item Exploración vs. explotación: balance entre probar y estabilizarse.
    \end{itemize}
  \end{block}
\end{frame}

\begin{frame}[t]{MDP: esquema de estados y acciones}
  \begin{center}
    \begin{tikzpicture}[
      state/.style={circle, draw, minimum size=1.05cm, thick, fill=gray!5},
      >={stealth},
      ->, thick,
      every node/.style={font=\small},
      scale=1.08,
      transform shape
    ]
      \node[state] (s1) at (0,1) {$s_1$};
      \node[state] (s2) at (2.4,2.2) {$s_2$};
      \node[state] (s3) at (2.4,-0.2) {$s_3$};
      \node[state] (s4) at (5,1) {$s_4$};
      
      \path (s1) edge node[above, sloped] {$a_1$} (s2);
      \path (s1) edge node[below, sloped] {$a_2$} (s3);
      \path (s2) edge node[above] {$a_3$} (s4);
      \path (s3) edge node[below] {$a_4$} (s4);
    \end{tikzpicture}
  \end{center}
\end{frame}

\begin{frame}[t]{MDP: Bellman en tiempo continuo}
  \begin{block}{Ecuación de Bellman}
    \[
      V(x) = \max_a \left\{ r(x,a)\Delta t + e^{-\beta\Delta t}\E[V(x')] \right\}
    \]

    En el límite $\Delta t\to 0$:
    \[
      \beta V(x) = \max_a \left\{ r(x,a) + \mathcal{A}_a V(x) \right\},
    \]
    donde $\mathcal{A}_a$ es el generador asociado a la acción $a$.
  \end{block}
\end{frame}

\begin{frame}[t]{Redes de Neuronas Espike (SNN)}
  \begin{columns}
    \column{0.5\textwidth}
    \begin{center}
      \begin{tikzpicture}[
        neuron/.style={circle, draw, minimum size=0.8cm},
        ->, >=stealth, thick
      ]
        \foreach \i in {1,2,3} {
          \node[neuron] (n\i) at (0,1.5*\i-2) {$n_\i$};
        }
        \node[neuron] (out) at (3,1) {$out$};
        
        \path (n1) edge node[above, sloped, font=\tiny] {$\alpha_{1}$} (out);
        \path (n2) edge node[above, font=\tiny] {$\alpha_{2}$} (out);
        \path (n3) edge node[below, sloped, font=\tiny] {$\alpha_{3}$} (out);
      \end{tikzpicture}
    \end{center}
    
    \begin{block}{Modelado con procesos de nacimiento}
      \begin{itemize}
        \item Potencial de membrana = estado
        \item Llegada de spikes = nacimientos
        \item Disparo = muerte (reset)
      \end{itemize}
    \end{block}
    
    \column{0.5\textwidth}
    \begin{block}{Tasas de disparo}
      \begin{align*}
        \lambda_n &= \text{tasa de entrada}\\
        \mu_n &= \text{tasa de disparo}\\
        \Prob\{\text{disparo}\} &= \frac{\mu_n}{\lambda_n+\mu_n}
      \end{align*}
      
      \vspace{0.3cm}
      \begin{center}
        \begin{tikzpicture}
          \draw[->, thick] (0,0) -- (4,0);
          \foreach \x in {0.3,0.8,1.5,2.2,3.1} {
            \draw[blue, thick] (\x,0) -- (\x,0.3);
          }
        \end{tikzpicture}
      \end{center}
    \end{block}
  \end{columns}
\end{frame}

\begin{frame}[t]{Sistemas de Recomendación}
  \begin{columns}
    \column{0.5\textwidth}
    \begin{center}
      \begin{tikzpicture}[
        item/.style={rectangle, draw, minimum width=1.2cm, minimum height=0.8cm},
        ->, >=stealth, thick
      ]
        \node[item] (i1) at (0,2) {Item 1};
        \node[item] (i2) at (2,2.5) {Item 2};
        \node[item] (i3) at (2,1.5) {Item 3};
        \node[item] (i4) at (4,2) {Item 4};
        
        \path (i1) edge node[above, sloped, font=\tiny] {$\alpha_{12}$} (i2);
        \path (i1) edge node[below, sloped, font=\tiny] {$\alpha_{13}$} (i3);
        \path (i2) edge node[above, sloped, font=\tiny] {$\alpha_{24}$} (i4);
        \path (i3) edge node[below, sloped, font=\tiny] {$\alpha_{34}$} (i4);
      \end{tikzpicture}
    \end{center}
    
    \begin{block}{Filtrado colaborativo}
      \begin{itemize}
        \item Usuarios = cadenas de Markov
        \item Transiciones = preferencias
        \item Recomendaciones = estados más probables
      \end{itemize}
    \end{block}
    
    \column{0.5\textwidth}
    \begin{block}{Tiempos medios de recomendación}
      \[ b(x) = \frac{1}{\alpha(x)} + \sum_{y\neq x,z} \frac{\alpha(x,y)}{\alpha(x)} b(y) \]
      
      \begin{center}
        \begin{tikzpicture}
          \begin{axis}[
            width=\textwidth,
            height=0.3\textheight,
            xlabel={ítem},
            ylabel={probabilidad},
            ybar,
            bar width=0.5cm
          ]
            \addplot coordinates {(1,0.3) (2,0.5) (3,0.2)};
          \end{axis}
        \end{tikzpicture}
      \end{center}
    \end{block}
  \end{columns}
\end{frame}

\section{Ecuaciones Forward y Backward}

\begin{frame}{Ecuaciones backward}
  Condicionando en el primer intervalo:
  \[ p_{t+\Delta t}(x,y) = \sum_z p_{\Delta t}(x,z) p_t(z,y). \]
  
  Usando probabilidades infinitesimales:
  \begin{align*}
    p_{t+\Delta t}(x,y) &= (1-\alpha(x)\Delta t)p_t(x,y) + \sum_{z\neq x} \alpha(x,z)\Delta t\, p_t(z,y) + o(\Delta t).
  \end{align*}
  
  Tomando $\Delta t\to 0$:
  \[ \frac{\partial}{\partial t} p_t(x,y) = -\alpha(x)p_t(x,y) + \sum_{z\neq x} \alpha(x,z)p_t(z,y). \]
  
  En forma matricial: $\displaystyle\frac{d}{dt}\PP_t = \A\PP_t$.
\end{frame}

\begin{frame}{Ecuaciones forward}
  Condicionando en el último intervalo:
  \[ p_{t+\Delta t}(x,y) = \sum_z p_t(x,z) p_{\Delta t}(z,y). \]
  
  Usando probabilidades infinitesimales:
  \begin{align*}
    p_{t+\Delta t}(x,y) &= p_t(x,y)(1-\alpha(y)\Delta t) + \sum_{z\neq y} p_t(x,z)\alpha(z,y)\Delta t + o(\Delta t).
  \end{align*}
  
  Tomando $\Delta t\to 0$:
  \[ \frac{\partial}{\partial t} p_t(x,y) = -\alpha(y)p_t(x,y) + \sum_{z\neq y} \alpha(z,y)p_t(x,z). \]
  
  En forma matricial: $\displaystyle\frac{d}{dt}\PP_t = \PP_t\A$.
\end{frame}

\begin{frame}{Comparación}
  \begin{center}
    \begin{tabular}{|c|c|}
      \hline
      \textbf{Backward} & \textbf{Forward} \\
      \hline
      $\frac{d}{dt}\PP_t = \A\PP_t$ & $\frac{d}{dt}\PP_t = \PP_t\A$ \\
      \hline
      Condiciona en $X_0$ & Condiciona en $X_t$ \\
      \hline
      Siempre válidas & Requieren regularidad \\
      \hline
    \end{tabular}
  \end{center}
  
  \vspace{0.5cm}
  
  Para espacios finitos, ambas tienen la misma solución:
  \[ \PP_t = e^{t\A}. \]
  
  \vspace{0.3cm}
  
  \begin{center}
    \begin{tikzpicture}
      \node at (0,0) {$e^{t\A}$};
      \draw[->, thick] (-1.5,0.5) -- (1.5,0.5) node[midway, above] {backward};
      \draw[->, thick] (-1.5,-0.5) -- (1.5,-0.5) node[midway, below] {forward};
    \end{tikzpicture}
  \end{center}
\end{frame}

\begin{frame}{Resumen visual}
  \begin{center}
    \begin{tikzpicture}[
      scale=0.64,
      transform shape,
      node distance=1.55cm,
      block/.style={rectangle, draw, minimum width=3.1cm, minimum height=0.9cm, rounded corners, thick, fill=gray!5, align=center}
    ]
      \node[block] (tasas) at (0,0) {Tasas $\alpha(x,y)$};
      \node[block] (gen) at (0,-1.7) {Generador $\A$};
      \node[block] (exp) at (0,-3.4) {$\PP_t = e^{t\A}$};
      \node[block] (est) at (4.9,-1.7) {Distribución\\estacionaria $\pi$};
      \node[block] (tiempos) at (-4.9,-1.7) {Tiempos medios\\de paso};
      \node[block] (bd) at (0,-5.15) {Nacimiento--muerte};

      \draw[->, thick] (tasas) -- (gen);
      \draw[->, thick] (gen) -- (exp);
      \draw[->, thick] (gen) -- (est);
      \draw[->, thick] (gen) -- (tiempos);
      \draw[->, thick] (exp) -- (bd);

      \node[draw, rounded corners, thick, fill=white, align=center, inner sep=4pt, text width=4.55cm] at (-5.9,-4.55) {
        $\pi\A = 0$ \\[3pt]
        $\alpha(x)b(x)=1+\sum_{y\neq x,z}\alpha(x,y)b(y)$
      };
    \end{tikzpicture}
  \end{center}
\end{frame}

\begin{frame}{Aplicaciones en IA - Resumen}
  \small
  \begin{center}
    \begin{tabular}{|p{0.34\textwidth}|p{0.46\textwidth}|}
      \hline
      \textbf{Aplicación IA} & \textbf{Concepto CTMC} \\
      \hline
      Modelos de lenguaje & Transiciones entre tokens y tiempos de atención \\
      Aprendizaje por refuerzo & MDP en tiempo continuo \\
      Redes SNN & Nacimiento y muerte de spikes \\
      Sistemas de recomendación & Tiempos medios de paso \\
      Optimización & Distribución estacionaria y estabilidad \\
      \hline
    \end{tabular}
  \end{center}
\end{frame}

\begin{frame}{Ventajas del modelado en tiempo continuo}
  \begin{block}{Ventajas principales}
    \begin{itemize}
      \item Captura dinámicas asíncronas.
      \item Modela eventos en tiempo real.
      \item Generaliza procesos discretos.
      \item Permite análisis de convergencia y estabilidad.
    \end{itemize}
  \end{block}
\end{frame}

\end{document}